\section*{Chapitre \ref{ch:g12:buffers-predominance} --- Solutions tampons et diagrammes de prédominance}

\begin{solution}{exo:g12:buffers-predominance:1}
\omterm{def:g9:acids-bases-ph:ph}{pH} 2~: \ce{CH3COOH}. \omterm{def:g9:acids-bases-ph:ph}{pH} 4.76~: les deux en quantités égales. \omterm{def:g9:acids-bases-ph:ph}{pH} 7~:
\ce{CH3COO-}.
\end{solution}

\begin{solution}{exo:g12:buffers-predominance:2}
\omterm{def:g9:acids-bases-ph:ph}{pH} 4~: $1/(1 + 10^{0.76}) = 0.15$. \omterm{def:g9:acids-bases-ph:ph}{pH} 6~: $1/(1 + 10^{-1.24}) = 0.95$.
\end{solution}

\begin{solution}{exo:g12:buffers-predominance:3}
Jaune~; vert (dans sa \omterm{def:g12:buffers-predominance:indicator}{zone de virage})~; bleu.
\end{solution}

\begin{solution}{exo:g12:buffers-predominance:4}
À \omterm{def:g9:acids-bases-ph:ph}{pH} 7.4, au-dessous de 9.26~: \ce{NH4+}. À \omterm{def:g9:acids-bases-ph:ph}{pH} 11~: \ce{NH3}.
\end{solution}

\begin{solution}{exo:g12:buffers-predominance:5}
Un indicateur est lui-même un \omterm{def:g12:ka-and-pka:bronsted}{couple acide/base}~: en grande quantité, il
modifierait le \omterm{def:g9:acids-bases-ph:ph}{pH} qu'il est censé indiquer.
\end{solution}

\begin{solution}{exo:g12:buffers-predominance:6}
$10^{\mathrm{pH} - 4.76}$~: 0.1~; 1~; 10.
\end{solution}

\begin{solution}{exo:g12:buffers-predominance:7}
La phénolphtaléine (8.0--10.0) ou le bleu de thymol (8.0--9.1)~: leurs
\omterm{def:g12:buffers-predominance:indicator}{zones de virage} contiennent 8.7. Le rouge de méthyle vire entre 4.2 et
6.3, loin de 8.7~: il aurait changé de couleur bien avant.
\end{solution}

\begin{solution}{exo:g12:buffers-predominance:8}
$\mathrm{pH} = 4.76 + \log(0.10/0.20) = 4.76 - 0.30 = 4.46$.
\end{solution}

\begin{solution}{exo:g12:buffers-predominance:9}
\omterm{def:g9:acids-bases-ph:ph}{pH} 1~: le \omterm{def:g9:ions:ion}{cation} \ce{H3N+-CH2-COOH}, de charge $+1$. \omterm{def:g9:acids-bases-ph:ph}{pH} 6~: le
zwitterion, de charge globale 0. \omterm{def:g9:acids-bases-ph:ph}{pH} 12~: l'\omterm{def:g9:ions:ion}{anion} \ce{H2N-CH2-COO-}, de
charge $-1$.
\end{solution}

\begin{solution}{exo:g12:buffers-predominance:10}
Eau~: de 7.0 à $-\log(\num{1.0e-4}/0.101) = 3.0$, soit une baisse de 4
unités.
Tampon~: de 4.76 à $4.76 + \log(9.9/10.1) = 4.75$, soit une baisse de
0.01.
\end{solution}

\begin{solution}{exo:g12:buffers-predominance:11}
\ce{NH4+}/\ce{NH3} ($pK_a = 9.26$). Rapport
$[\ce{NH3}]/[\ce{NH4+}] = 10^{9.5 - 9.26} = 1.7$.
\end{solution}

\begin{solution}{exo:g12:buffers-predominance:12}
Le \omterm{def:g9:acids-bases-ph:ph}{pH} a diminué d'une unité quand le rapport atteint $1/10$~:
$(10 - n)/(10 + n) = 0.1$, donc $n = \qty{8.2}{mmol}$. Un tampon qui
contient davantage de chaque forme peut absorber plus d'acide ou de base
pour une même variation du rapport~: son pouvoir tampon est plus grand.
\end{solution}

\begin{solution}{exo:g12:buffers-predominance:13}
Rapport $10^{5.0 - 4.76} = 1.74$~; l'acide représente \qty{10}{mmol},
il faut donc
\qty{17.4}{mmol} d'éthanoate, soit \qty{174}{mL} de la solution.
\end{solution}

\begin{solution}{exo:g12:buffers-predominance:14}
Il possède deux \omterm{def:g12:ka-and-pka:bronsted}{couples acide/base}, avec deux $pK_a$, d'où deux
changements de couleur~: rouge à \omterm{def:g9:acids-bases-ph:ph}{pH} 1, jaune à \omterm{def:g9:acids-bases-ph:ph}{pH} 5, bleu à \omterm{def:g9:acids-bases-ph:ph}{pH} 10.
\end{solution}

\begin{solution}{exo:g12:buffers-predominance:15}
La \omterm{def:g10:concentration-and-dilution:dilution}{dilution} divise les deux concentrations par 10~: le rapport, et donc
le \omterm{def:g9:acids-bases-ph:ph}{pH} (4.76), ne changent pas. Un \omterm{def:g12:ka-and-pka:strong-acid}{acide fort} dilué dix fois voit son \omterm{def:g9:acids-bases-ph:ph}{pH}
augmenter d'une unité.
\end{solution}

\begin{solution}{pb:g12:buffers-predominance:1}
\textbf{1.} \ce{CO2 + 2H2O <=> HCO3- + H3O+}.

\textbf{2.} Le dioxyde de carbone dissous (avec l'eau) est l'acide,
l'\omterm{def:g9:ions:ion}{ion} hydrogénocarbonate la base.

\textbf{3.} La valeur de la table correspond à \qty{25}{\celsius} dans
l'eau pure~; le sang est à \qty{37}{\celsius} et contient beaucoup
d'autres \omterm{def:g9:ions:ion}{ions}, qui modifient la constante (et les physiologistes
comptent tout le dioxyde de carbone dissous).

\textbf{4.} À \qty{24}{mmol/L}, $0.024 \times 5.0 = \qty{0.12}{mol}$.

\textbf{5.} Un axe de \omterm{def:g9:acids-bases-ph:ph}{pH} coupé à 6.1~: \ce{CO2} au-dessous, \ce{HCO3-}
au-dessus.

\textbf{6.} À 7.4, au-dessus de 6.1~: l'hydrogénocarbonate.

\textbf{7.} Légèrement basique~: son \omterm{def:g9:acids-bases-ph:ph}{pH} est supérieur à 7.

\textbf{8.} $10^{7.4 - 6.1} = 10^{1.3} = 20$.

\textbf{9.} $24 / 20 = \qty{1.2}{mmol/L}$.

\textbf{10.} $10^{1.28} = 19$ et $10^{1.32} = 21$.

\textbf{11.} $20/21 = \qty{95}{\percent}$.

\textbf{12.} $10^{7.6 - 6.1} = 10^{1.5} = 32$.

\textbf{13.} L'\omterm{def:g9:ions:ion}{ion} hydrogénocarbonate~:
\ce{HCO3- + H3O+ -> CO2 + 2H2O}.

\textbf{14.} $10^{7.1 - 6.1} = 10$.

\textbf{15.} De $24/20 = 1.2$ à $24/10 = \qty{2.4}{mmol/L}$~: il double.
Les poumons respirent plus vite et plus profondément, pour chasser
l'excès de dioxyde de carbone.

\textbf{16.} Un \omterm{def:g12:ka-and-pka:strong-acid}{acide faible} et sa base peuvent chacun absorber ce que
l'on ajoute, et le \omterm{def:g9:acids-bases-ph:ph}{pH} ne dépend que de leur rapport~; un \omterm{def:g12:ka-and-pka:strong-acid}{acide fort} ne
ferait que pousser le \omterm{def:g9:acids-bases-ph:ph}{pH} dans un seul sens.

\textbf{17.} Pas du tout~: le rapport, et donc le \omterm{def:g9:acids-bases-ph:ph}{pH}, restent les mêmes.

\textbf{18.} Environ 20.
\end{solution}
